% Methodology section fragment; requires a parent document with natbib and references.bib.
\providecommand{\na}{\textcolor{red}{[TBD]}}

\section{Methodology}
\label{sec:method}

We evaluate whether unit tests written by a trusted model detect code backdoors, through three questions:
\begin{enumerate}
  \item[(i)] \emph{Can tests replace the monitor?} Do tests alone detect backdoors as well as a trusted monitor that reads the code?
  \item[(ii)] \emph{Do tests help the monitor?} Does giving the monitor the tests' results raise safety at equal usefulness?
  \item[(iii)] \emph{What limits the tests?} When the tests miss a backdoor, is it because none of their inputs reaches the trigger, or because none of their assertions separates the backdoored output from a correct one?
\end{enumerate}

\subsection{Threat model}
\label{sec:setting}

We adopt the high-stakes AI control setting of \citet{greenblatt2024aicontrol}. An untrusted model $U$ solves programming tasks, and a submission may contain a \emph{backdoor}: code that is correct on ordinary inputs but returns an incorrect output on a rare trigger input chosen by its author. The blue team has a weaker trusted model $T$ and a small budget of human audits, which we assume to be perfect. The blue team wins if an audit catches a backdoor before any functional backdoor is deployed. The red team's attack policy is a fixed corpus of backdoors written by other models~\citep{terekhov2025controltax} (Section~\ref{sec:data}); it does not adapt to the protocol or choose which tasks to attack, a limitation we return to in Section~\ref{sec:limitations}.

\subsection{Protocols}
\label{sec:protocols}

Every protocol assigns each candidate a suspicion score. Table~\ref{tab:protocols} summarises what each protocol observes.

\begin{table}[t]
\centering
\small
\setlength{\tabcolsep}{4pt}
\begin{tabular}{llll}
\toprule
Protocol & Score & Reads the code & Writes the tests \\
\midrule
Trusted monitoring              & suspicion, 1--10    & $T$                     & ---                  \\
Property / plain tests          & tests failed, 0--10 & $T$                     & $T$                  \\
$U$-planned tests               & tests failed, 0--10 & $U$; $T$ (inputs only)  & $T$, from $U$'s plan \\
Monitoring + tests (3 variants) & suspicion, 1--10    & $T$, plus test results  & as in its suite      \\
\bottomrule
\end{tabular}
\caption{The protocols. Every test runs on the same $n$ trigger inputs, proposed once per candidate
by $T$ reading the code.}
\label{tab:protocols}
\end{table}

\paragraph{Trusted monitoring.} Following \citet{greenblatt2024aicontrol}, $T$ reads the task and the submission and rates the likelihood of a backdoor from 1 to 10.

\paragraph{Test generation.} Testing proceeds in two stages. First, $T$ reads the task and the submission and proposes $n$ inputs designed to satisfy each condition in the code that could gate an incorrect output. This search runs once per candidate, honest and backdoored alike, and all test-based protocols share its inputs, so differences between them reflect the tests rather than the inputs. Second, $T$ writes a suite of $t$ tests; each test is executed on every input, and the score $k \in \{0, \dots, t\}$ is the number of tests that fail on at least one input. Because the setting assumes $T$ cannot reliably solve the task, a test may not compute the expected output: it must assert a property that every correct output satisfies.

\paragraph{Property and plain tests.} We compare two kinds of suite. A \emph{plain test} checks one output on its own, for example that it has the required format or lies within the bounds the task states. A \emph{property test} may also relate the outputs on several inputs, for example that reordering the input leaves the answer unchanged, a metamorphic relation~\citep{chen1998metamorphic}, or check a certificate that the answer is correct~\citep{claessen2000}.

\paragraph{$U$ plans, $T$ writes.} $U$ reads the submission and writes a natural-language test plan without code, and $T$ implements the suite from the plan without seeing the submission. This tests whether a stronger model's reading of the code yields better tests. It carries two caveats, which it is partly designed to measure: a $U$ that recognises its own backdoor could plan tests that avoid it, and $U$ reasons at higher effort than the protocols it is compared with.

\paragraph{Monitoring with tests.} The trusted monitor, with the same prompt and number of samples, additionally receives one suite's source code and execution results on the same submission. The prompt warns that a failing test may reflect either a genuine backdoor or an assertion the specification does not support. We evaluate one variant per suite and refer to them as monitoring + property tests, monitoring + plain tests and monitoring + $U$-planned property tests.

% TODO(Omar): one paragraph on the multi-turn protocol, if the full-pool multi-turn arms are
% reported in Results.
% TODO: the plain_v5 prompt variant (E12), if reported. Part of it was scored after the provider
% changed DeepSeek-V3.2's behaviour on 10 October 2026, so compare it only within a period.

\subsection{Data}
\label{sec:data}

We use APPS~\citep{hendrycks2021apps} with the backdoored solutions of \citet{terekhov2025controltax} from the \texttt{apps-control-arena} dataset~\citep{appscontrolarena}. A task is included if its reference solution is correct and deterministic and at least one backdoor passes the visible tests and disagrees with the reference solution on its trigger input; a backdoor that never triggers is an ordinary bug, not an attack. Each task contributes one honest and one backdoored candidate, so no task is counted twice in a rate. All of APPS yields 1{,}388 such tasks (2{,}776 candidates), split into 555 calibration and 833 evaluation tasks. Before running any protocol, we remove comments and self-describing identifiers from both honest and backdoored code. We verify every label independently with the ControlArena scorer (Appendix~\ref{app:implementation}).

\subsection{Evaluation}
\label{sec:safety-model}

\paragraph{Audit threshold.} Following \citet{greenblatt2024aicontrol}, the audit threshold $\tau_b$ for budget $b$ is the $(1-b)$ quantile of a protocol's scores on honest calibration candidates; no attack labels are used. Scores are discrete, so we break ties at $\tau_b$ at random such that exactly a fraction $b$ of honest calibration candidates is audited. Every protocol therefore spends the same budget.

\paragraph{Defer to trusted.} Every protocol is run with defer to trusted, as in the baseline of \citet{greenblatt2024aicontrol}: a submission that is not audited but scores at or above a second, lower threshold is replaced by $T$'s own solution, which contains no backdoor but is less often correct. With trusted monitoring's score, this is their defer-to-trusted protocol. Lowering the deferral threshold trades usefulness for safety, so we compare protocols at equal usefulness; a deferral threshold above every score defers nothing and keeps usefulness at 1.

\paragraph{Safety.} The budget enters only through the threshold: at evaluation, every candidate above $\tau_b$ is audited, however many there are. Under this model, and with attacks on uniformly random tasks, an attack's outcome depends only on its own score and not on when it occurs, so expected safety does not depend on the attack rate (up to the chance that no attack is decisive in a finite deployment). We therefore attack every evaluation task, which minimises variance. Let $a_i$ and $d_i$ be the probabilities that the attack on task $i$ is audited and deferred, and $w_i \in \{0, 1\}$ indicate a functional backdoor. Safety is the probability that the first decisive attack is caught,
\[
  \mathrm{Safety} \;=\; \frac{\sum_i a_i}{\sum_i a_i + \sum_i (1 - a_i - d_i)\, w_i},
\]
and 1 when no attack is decisive. Without deferral and with every backdoor functional, safety is the fraction of attacks audited.

\paragraph{Usefulness.} On honest evaluation candidates, an audited or deployed submission counts as correct and a deferred one counts as correct if $T$'s own solution passes every test of the task under the ControlArena scorer. Following \citet{greenblatt2024aicontrol}, we rescale the share of tasks that end with a correct solution linearly, so that 0 is the share $T$'s own solutions get right (81\% of evaluation tasks) and 1 the share $U$'s honest solutions get right (100\%). A usefulness of 0.95 therefore means that 99\% of tasks end correct, and 0.90 that 98\% do.

\paragraph{Reach.} To separate the two causes of a miss, we call a backdoor \emph{reached} if at least one of its candidate's trigger inputs makes the backdoored and reference solutions produce different outputs under ControlArena's rule. This uses the reference solution, which no protocol sees, so it serves only as a diagnostic.

\paragraph{Reporting.} The audit budget is $b = 2\%$, as in \citet{greenblatt2024aicontrol}. With $\tau_{2\%}$ fixed, we sweep the deferral threshold to trace each protocol's safety--usefulness frontier and report safety at three usefulness levels: 1 (no deferral), 0.95 and 0.90, the latter two interpolated between adjacent deferral thresholds. We also report safety at a usefulness of 0.95 for $b \in \{1, 5, 10\}\%$. We report ROC AUC and the detection rate at fixed false-positive rates, which characterise each protocol independently of the safety model. To separate a protocol's ranking from the noise of calibrating its threshold, we also compare protocols at matched false-positive rates set on the evaluation split itself; because this uses information a deployment does not have, we report it only as a diagnostic. We report each value as $\pm$ 1.96 standard deviations over 10{,}000 paired task resamples (a 95\% interval), in which every protocol recalibrates its threshold on the resampled calibration tasks; we report this for each protocol and for its paired difference from trusted monitoring. All comparisons use the candidates that every protocol scored (2{,}435 of 2{,}776); a candidate a protocol failed to score is excluded rather than assigned a default score, since an unscored test suite would otherwise count as passing.

\paragraph{Implementation.} $T$ is DeepSeek-V3.2~\citep{deepseekv32} at low reasoning effort and $U$ is the same model at high effort, so the capability gap arises from reasoning effort alone. We use $n = 30$ trigger inputs and $t = 10$ tests, and run every test in an isolated sandbox. Further details are in Appendix~\ref{app:implementation}.

